%% ============================================================
%%  PLANTILLA DE APUNTES DE MATEMÁTICA
%%  Compilar con:  pdflatex plantilla_apuntes.tex   (2 veces)
%%  Para el índice analítico: makeindex plantilla_apuntes
%%  y luego pdflatex otra vez. (Con --shell-escape se hace solo.)
%% ============================================================
\documentclass[11pt,a4paper,oneside]{report}

%% ---------- Idioma y codificación ----------
\usepackage[T1]{fontenc}
\usepackage[utf8]{inputenc}
\usepackage{lmodern}
\usepackage[spanish,es-noshorthands]{babel}
\usepackage{microtype}

%% ---------- Página y párrafos ----------
\usepackage[margin=2.5cm]{geometry}
\usepackage{parskip}            % separación entre párrafos sin sangría
\usepackage{setspace}
\setstretch{1.1}
\usepackage{fancyhdr}
\pagestyle{fancy}
\fancyhf{}
\fancyhead[L]{\small\nouppercase{\leftmark}}
\fancyhead[R]{\small\thepage}
\renewcommand{\headrulewidth}{0.4pt}

%% ---------- Matemática ----------
\usepackage{amsmath,amssymb,amsthm,mathtools}
\usepackage{bm}

%% ---------- Figuras, tablas, colores ----------
\usepackage{graphicx}
\graphicspath{{figuras/}}       % pon tus imágenes en la carpeta "figuras"
\usepackage{xcolor}
\usepackage{booktabs}           % tablas bonitas (\toprule, \midrule)
\usepackage{array}
\usepackage{tabularx}
\usepackage{longtable}
\usepackage{caption}
\usepackage{float}              % permite [H] para fijar posición
\usepackage{mwe}                % solo para imágenes de ejemplo; bórralo luego
\usepackage{tikz}

%% ---------- Código ----------
\usepackage{listings}
\lstset{
  basicstyle=\ttfamily\small,
  keywordstyle=\color{blue!70!black}\bfseries,
  commentstyle=\color{green!40!black}\itshape,
  stringstyle=\color{orange!80!black},
  numbers=left, numberstyle=\tiny\color{gray}, numbersep=8pt,
  backgroundcolor=\color{gray!8},
  frame=single, rulecolor=\color{gray!40},
  breaklines=true, showstringspaces=false, tabsize=4,
  literate={á}{{\'a}}1 {é}{{\'e}}1 {í}{{\'i}}1 {ó}{{\'o}}1 {ú}{{\'u}}1
           {Á}{{\'A}}1 {É}{{\'E}}1 {Í}{{\'I}}1 {Ó}{{\'O}}1 {Ú}{{\'U}}1
           {ñ}{{\~n}}1 {Ñ}{{\~N}}1
}

%% ---------- Índices ----------
\usepackage{imakeidx}
\makeindex[title=Índice analítico, intoc]

%% ---------- Cajas de colores ----------
\usepackage[most]{tcolorbox}

%% ---------- Hipervínculos y referencias (SIEMPRE al final) ----------
\usepackage[colorlinks=true,linkcolor=blue!60!black,citecolor=green!40!black,
            urlcolor=blue!60!black]{hyperref}
\usepackage[spanish,nameinlink]{cleveref}

%% ============================================================
%%  ENTORNOS: definición, teorema, demostración, notas...
%%  Todos comparten UN solo contador, dentro de cada capítulo.
%% ============================================================
\theoremstyle{definition}
\newtheorem{definicion}{Definición}[chapter]
\newtheorem{ejemplo}[definicion]{Ejemplo}
\newtheorem{ejercicio}[definicion]{Ejercicio}

\theoremstyle{plain}
\newtheorem{teorema}[definicion]{Teorema}
\newtheorem{proposicion}[definicion]{Proposición}
\newtheorem{lema}[definicion]{Lema}
\newtheorem{corolario}[definicion]{Corolario}

\theoremstyle{remark}
\newtheorem*{observacion}{Observación}
\newtheorem*{notademo}{Nota sobre la demostración}
\newtheorem*{idea}{Idea clave}

\renewcommand{\proofname}{Demostración}
\renewcommand{\qedsymbol}{$\blacksquare$}

% Nombres para cleveref
\crefname{definicion}{definición}{definiciones}
\Crefname{definicion}{Definición}{Definiciones}
\crefname{teorema}{teorema}{teoremas}
\Crefname{teorema}{Teorema}{Teoremas}
\crefname{proposicion}{proposición}{proposiciones}
\Crefname{proposicion}{Proposición}{Proposiciones}
\crefname{lema}{lema}{lemas}
\Crefname{lema}{Lema}{Lemas}
\crefname{corolario}{corolario}{corolarios}
\Crefname{corolario}{Corolario}{Corolarios}
\crefname{ejemplo}{ejemplo}{ejemplos}
\Crefname{ejemplo}{Ejemplo}{Ejemplos}
\crefname{ejercicio}{ejercicio}{ejercicios}
\Crefname{ejercicio}{Ejercicio}{Ejercicios}
\crefname{equation}{ecuación}{ecuaciones}
\Crefname{equation}{Ecuación}{Ecuaciones}
\crefname{figure}{figura}{figuras}
\Crefname{figure}{Figura}{Figuras}
\crefname{table}{tabla}{tablas}
\Crefname{table}{Tabla}{Tablas}
\crefname{chapter}{capítulo}{capítulos}
\Crefname{chapter}{Capítulo}{Capítulos}
\crefname{section}{sección}{secciones}
\Crefname{section}{Sección}{Secciones}

% Estilo visual: barra lateral de color + fondo suave
\tcbset{
  cajabase/.style={enhanced, breakable, sharp corners, boxrule=0pt,
    frame hidden, left=8pt, right=6pt, top=4pt, bottom=4pt,
    before skip=10pt, after skip=10pt}
}
\tcolorboxenvironment{definicion}{cajabase,
  colback=blue!6,   borderline west={3pt}{0pt}{blue!65!black}}
\tcolorboxenvironment{teorema}{cajabase,
  colback=red!6,    borderline west={3pt}{0pt}{red!70!black}}
\tcolorboxenvironment{proposicion}{cajabase,
  colback=orange!7, borderline west={3pt}{0pt}{orange!80!black}}
\tcolorboxenvironment{lema}{cajabase,
  colback=orange!7, borderline west={3pt}{0pt}{orange!80!black}}
\tcolorboxenvironment{corolario}{cajabase,
  colback=orange!7, borderline west={3pt}{0pt}{orange!80!black}}
\tcolorboxenvironment{ejemplo}{cajabase,
  colback=green!5,  borderline west={3pt}{0pt}{green!45!black}}
\tcolorboxenvironment{ejercicio}{cajabase,
  colback=green!5,  borderline west={3pt}{0pt}{green!45!black}}
\tcolorboxenvironment{notademo}{cajabase,
  colback=yellow!14, borderline west={3pt}{0pt}{yellow!60!black}}
\tcolorboxenvironment{observacion}{cajabase,
  colback=gray!10,  borderline west={3pt}{0pt}{gray!70}}
\tcolorboxenvironment{idea}{cajabase,
  colback=purple!6, borderline west={3pt}{0pt}{purple!70}}

%% ============================================================
%%  ATAJOS MATEMÁTICOS
%% ============================================================
\newcommand{\N}{\mathbb{N}}
\newcommand{\Z}{\mathbb{Z}}
\newcommand{\Q}{\mathbb{Q}}
\newcommand{\R}{\mathbb{R}}
\newcommand{\C}{\mathbb{C}}
\newcommand{\eps}{\varepsilon}
\newcommand{\dd}{\mathrm{d}}
\newcommand{\dpar}[2]{\frac{\partial #1}{\partial #2}}
\DeclarePairedDelimiter{\abs}{\lvert}{\rvert}
\DeclarePairedDelimiter{\norm}{\lVert}{\rVert}
\DeclarePairedDelimiter{\pesc}{\langle}{\rangle}
\DeclarePairedDelimiter{\conj}{\{}{\}}
\DeclareMathOperator{\sen}{sen}
\DeclareMathOperator{\tg}{tg}
\DeclareMathOperator{\Img}{Im}
\DeclareMathOperator{\Ker}{Ker}
\DeclareMathOperator{\rg}{rg}

% Atajo para "término indexado": lo escribe en el texto Y lo añade al índice
\newcommand{\termino}[1]{\textbf{#1}\index{#1}}

%% ============================================================
%%  PORTADA
%% ============================================================
\title{Apuntes de Programación}
\author{Práxedis Jiménez Ruvalcaba}
\date{\today}

\begin{document}

\maketitle
\tableofcontents

\part{Aclaraciones}
Estos son los apuntes más particulares y cómo tal, es necesario tener una serie de aclaraciones para
entender la manera en que se hicieron y que se tiene pensado su uso. Para empezar, todo de aquí ha venido
de manera autodidácta y de una variedad más que diversa de fuentes. Mientras el resto de apuntes vienen
de un par de libros y de los propios cursos, aquí vienen de eso mismo, así como muchísimos foros, vídeos,
conversaciones, publicaciones, webs y de una cantidad irrasteable de fuentes, así que hay que tener precaución.\\

Lo segundo es que mucho, sino que todo, viene de preguntas que se concibieron durante la realización del código. 
De eso vienen siendo muchos proyectos, muchos con repositorios propios y algunos sin publicarse y otros que existen 
en los compendios donde tengo montones de ejercicios de práctica. Estos repositorios vienen siendo el compilado 
de ejercicios de C++ \url{https://github.com/PraxedisJRuv/practice},un framework de simulaciones en C++ 
\url{https://github.com/PraxedisJRuv/simulations}, pipline de IA para investigación clínica
\url{https://github.com/PraxedisJRuv/Health_Research}, compilado de ejercicos para python 
\url{https://github.com/PraxedisJRuv/Python-practice}. En general, todo lo que tiene entre mis repositorios en 
\url{https://github.com/PraxedisJRuv}.\\

Lo tercero es que no hay una seguridad clara de que estos apuntes sean para compartir más que a nivel 
personal para juntar apuntes desorganizados de alguna forma y para tener algo de lo que estudiar con referencias 
en caso de querer repasar o retomar algo. No es como el resto de apuntes que pueden servir como material 
complementario para el estudio de un curso o incluso apoyo para el diseño de uno. Este está poco recomendado para
ello. Para construir algo, programar algo, empaparse de lenguaje, herramientas y conceptos que no ve uno de primeras
ni se revisan en cursos de programación, sí sirve, esa es una buena idea. Unos apuntes que tener al lado cuando se 
construye algo.\\

Este es el compendio de apuntes con más $spanglish$ con diferencia, parte de la razón tiene que 
ver que el lenguaje usual en el ambiente del cómputo es en inglés y hay muchas palabras que no son 
cómodamente traducibles al español además de que es rarísimo encontrarles así y es más cómodo en el
propio inglés. De hecho, había apuntes que directamente mezcalaba no solo en los términos pero en las
propias oraciones y estrucutras al inglés con el español. Si intentó corregir estos detalles al transcribir 
los apuntes, pero siguen habiendo detalles. Así mismo, aquí se ignora la regla de cursivas para términos 
en lengua extranjera ya que son demasiados y sería algo asburdo.\\

\part{Generalidades}

\chapter{Vocabulario de Sistemas}

\section{Bases}
\begin{definicion}[\termino{Threads}]
  Unidades que realizan tareas separadas al mismo tiempo. Los threads del mismo 
  proceso comparten memoria y archivos, ergo, comunicarse rápido.\\
  Cada thread mantiene seguimiento de su ejecución incluido program counter, register set y 
  espacio en el stack para variables locales.
\end{definicion}
Hay en software/CPU como uno descargando y otro respondiendo al cursos y otro al buscador.\\
Hay CPU/Hardware que sirve para mucho software al mismo tiempo. \\
Multithreading funcionan porque las computadores tienen múltiples núcleos.

\section{Lock-Free}

\begin{definicion}[\termino{Non-Blocking Progress}]
  Se apoya en primitivos del hardware (compare-and-swap) CAS o 
  Fetch-And-Add (FAA) a través de operaciones atómicas garantizando que al menos 
  un hilo porque globalmente (look-free) o que todo hilo termine con una cantidad acotada de 
  pasos (wait-free).\\
  En general es la idea de diseño de organizar todo el flujo de manera apropiada para que no estorben los 
  elementos propios del sistema, que no se bloqueen entre si.
\end{definicion}

\begin{definicion}[\termino{ABA problem}]
  Cuando la dirección de memoria cambia de A a B y de vuelta a A,
  baitea o engaña una aparición CAS a que no pase nada, que se resuelve con poderosos versionados
  (Atomic Stomped Reference).
\end{definicion}

\begin{definicion}[\termino{Safe Memory Reclamation (SMR)}]
  Descolocar Nodos de manera segura sin usar after-free requiere de herramientas 
  como hazard pointers, epoch-based relocation, o garbage collection, más que 
  solo liberar celdad o espacios de memoria.
\end{definicion}

\section{Cache-Aware}

\begin{definicion}[\termino{False Sharing}]
  Cuando hilos independientes modifican variables separadas que residen en el mismo
  bloque de memoria, misma cache line de 64-byte CPU, causando problemas serios en la cache y memoria.
  Se mitiga alineando los bloques de memoria apropiadamente. 
\end{definicion}

\begin{definicion}[\termino{Locality and Prefetching}]
  Se diseñan explícitamente toda la estructura para maximizar la localización temporal 
  y especial, manteniendo nodos frecuemente atendidos con metadata "ho" en ciertas caches,
  reduciendo los viajes de memoria.
\end{definicion}

\begin{definicion}[\termino{Lazy Updates}]
  Difiere estructuras de sincornización pesadas en lotes, lo que 
  disminuye la tarifa de interconección en modelos de memoria baja. 
\end{definicion}
Es como el truco de juntar varias request de una API para no tener que estar 
evaluando cada proceso separado, mejor juntarlos en lotes y después evaluarlos todos
juntos sin tener que estar reiniciando la máquina todo el tiempo, aunque temrina quedando mucho
más sensible a errores ya que si una request del lote falla, rompería ese ciclo de optimización.

\section{Sistemas Distribuidos}

\begin{definicion}[\termino{Nodo}]
  Dispositivo de hardware o proceso de software que forma
  parte del sistema y se comunica por una red.
\end{definicion}

\begin{definicion}[\termino{Escalabilidad horizontal}]
  Capacidad de añadir más máquinas para manejar mayor carga de trabajo en
  lugar de manejar un solo equipo.
\end{definicion}

\begin{definicion}[\termino{Tolerancia a fallos}]
  Habilidad de la red para seguir funcionando aunque falle uno o más componentes.
\end{definicion}

\begin{definicion}[\termino{Teorema CAP}]
  Un sistema distribuidosolo puede garantizar 2 propiedades simultáneamente entre consistencia, 
  disponibilidad y tolerancia a particiones.
\end{definicion}
Esto debería ser un teorema, pero no tenía de dónde sacar la referencia para expresarlo
formalmente. 

\begin{definicion}[\termino{Consenso}]
  Proceso mediante el cual varios nodos se ponen de acuerdo sobre un mismo valor o estado
  del sistema.
\end{definicion}

\begin{definicion}[\termino{Idempotencia}]
  Propiedad de una operación de ejecutarla varias veces obteniendo el mismo 
  resultado y sin causar efectos secundarios no deseados.
\end{definicion}

\begin{definicion}[\termino{Balanceador de carga}]
  Componente que represeta el tráfico de peticiones entre mútilpes
  servidores para evitar saturaciones.
\end{definicion}
%% ============================================================
\part{C++}
%% ============================================================

\chapter{Bases}
Casi todos estos apuntes han sido obtenidos por partes y realizados por mi cuenta. A diferencia de los apuntes de 
matemáticas, muchos no van a ser expresamente formales, ya que no tengo dicha formación en ciencias computacionales,
al menos al momento de hacer los apuntes. Esto también es algo separado de mis proyectos de código donde a su vez hay
documentación, pero creo este es un posible espacio apropiado para revisar y exponer mucho de lo visto y entendido en
dichos proyectos. Solo quiero remarcar que las definiciones no son propiamente formales y probablemente falte detalle
y propiedad que habría en los de matemáticas.
\section{Punteros}


\begin{definicion}[Stack]\label{def:stack}
  .\\
  \begin{itemize}
    \item Sigue LIFO estricto
    \item Súper rápida pero con poca capacidad
    \item Sirve para parámetros o variables auxiliares
  \end{itemize}
\end{definicion}

\begin{definicion}[Heap]\label{def:heap}
  .\\
  \begin{itemize}
    \item Memoria dinámica para cuando se desarmen los datos o cuando tienen que vivir más allá de una memoria
    \item No sigue LIFO necesariamente
    \item Es la memoria que debe llamarse
  \end{itemize}
\end{definicion}

\begin{lstlisting}[language=C++, caption={new}]
// Reserva memoria para "size" enteros y retorna la dirección de memoria del primer valor
new int[size];
\end{lstlisting}

\begin{lstlisting}[language=C++, caption={const pointer}]
/*
No puede modificar  *ptr pero sí la dirección, es decir, 
no puede modificar x (valor) pero sí la dirección del puntero
*/
const int* ptr = &x;
\end{lstlisting}

\begin{lstlisting}[language=C++, caption={pointer to pointer}]
// Puntero hacía un puntero int*
int**
\end{lstlisting}
Para modificar la dirección a la que un puntero apunta dentro de una función, es necesario pasar la
referencia del puntero, que es un puntero de un puntero. \\

Cuando dos punteros apuntan al mismo bloque de memoria, si se restan, retornan la cantidad de elementos entre ellos,
no de bytes.


\printindex

\end{document}
